\chapter{Hip-Hop: A World Inside One Loop}
\label{ch:genre-hiphop-deep}

\classquote{凡音之起，由人心生也。人心之动，物使之然也。}{Book of Rites, ``Record of Music''}{Sound arises from the heart; when the heart is moved by things, it takes shape as sound.}

\section{What Hip-Hop Is: Rap Is Only the Front}
\label{sec:hiphop-what}

Many listeners meet hip-hop as one person talking very fast. That is only the front of the stage. Hip-hop is a culture that grew out of community rooms: a DJ extends the drum break of a record with two turntables and a mixer; an MC rhymes over that break; a B-boy or B-girl dances to it; a writer paints a name on a wall. Rapping is the line that records and travels most easily, but it is not the whole culture.

Writing a hip-hop song is therefore not ``reading the lyrics faster.'' It works on four things at once:

\begin{longtable}{p{0.16\textwidth}p{0.34\textwidth}p{0.38\textwidth}}
\toprule
What the ear notices & What the listener hears first & The writing question \\
\midrule
\endfirsthead
\toprule
What the ear notices & What the listener hears first & The writing question \\
\midrule
\endhead
Groove & The body nods before the mind understands the words & Which beat is the pocket? Where does the vocal sit against the drums? \\
Space & The loop breathes; the low end leaves room & If the sample is already full, should a pad be added at all? \\
Story & A specific person, place, or moment is audible & Is this verse telling one thing, or only rhyming? \\
Hook & A short phrase, sound, or gesture can be repeated & What is the smallest memorable unit? Does the song still work without it? \\
\bottomrule
\end{longtable}

The danger in hip-hop is not roughness; it is a lack of choices. If every beat is full, every line is at the same speed, and every layer plays at once, the listener has nowhere to enter. Craft often shows up as subtraction: half a beat of silence after a line so the next line has a landing place; one gap in a loop so the drums and 808 have somewhere to speak.

\section{The Route: From a Community Room to Global Playlists}
\label{sec:hiphop-history}

\subsection{The 1970s: The Moment the Break Was Extended}
\label{sec:hiphop-history-1970s}

On August 11, 1973, at 1520 Sedgwick Avenue in the Bronx, Cindy Campbell threw a back-to-school party to raise money. Her brother Clive Campbell---later DJ Kool Herc---used two turntables and a mixer to connect and repeat the hottest drum breaks in funk and soul records. Herc's friend Coke La Rock shouted names and short rhymes between the breaks to lift the room.{\textsuperscript{[1]}}

This is not the myth of one person inventing a genre; it is a checkable technical starting point. The breakbeat---isolating and extending the part of a record made for dancing---gave dancers more body time, and the MC's short phrases grew inside that time. Sampling, the groove pattern, and call-and-response all have roots here.

\subsection{The 1980s--1990s: From the Street to the Studio}
\label{sec:hiphop-history-1980s}

In the 1980s, rap moved from live spaces into the record industry. In 1982, Afrika Bambaataa's ``Planet Rock'' brought electronic timbres into hip-hop; Grandmaster Flash and the Furious Five's ``The Message'' put narrative rap on the map; Run-DMC used harder drums, shorter lines, and street fashion to make a version that the mainstream could see. In 1988, Public Enemy layered samples, noise, and political speech, while N.W.A pushed West Coast hardcore narrative to an extreme.

The 1990s are often called the golden age. New York boom-bap chopped samples more finely and cared intensely about drum placement; West Coast G-funk used slower grooves, fatter synthesizers, and funk samples; Nas, the Notorious B.I.G., Wu-Tang Clan, and Jay-Z made rap albums into complete listening works.{\textsuperscript{[2]}} The most important inheritance from this period is not a timbre but an idea: a beat can be designed like a composition.

\subsection{The 2000s--2020s: Low End, Trap, and Local Accents}
\label{sec:hiphop-history-2000s}

In the 2000s, Southern crunk, snap, and trap pushed the low end and the drum machine to the center. T.I.'s 2003 album \emph{Trap Muzik} put the word ``trap'' into the mainstream; producers such as Lex Luger, 808 Mafia, and Metro Boomin turned long-decay 808s, continuous hi-hats, and triplet flows into a reproducible language.{\textsuperscript{[3]}} In the 2010s, trap spread from Atlanta into pop, electronic music, and rap worldwide; Kendrick Lamar, Drake, and Travis Scott pulled rap apart into melody, atmosphere, stage, and visual event.

In the 2020s, drill traveled from Chicago's South Side to London and back into global short video; streaming and algorithms let one song borrow trap, afrobeat, hyperpop, and local folk languages at the same time. Genre borders are looser, but the arrangement questions have not changed: how does a loop keep space, a voice keep a pocket, and a section keep a hook?

\subsection{Chinese Hip-Hop: From an Attempt to an Accent}
\label{sec:hiphop-history-china}

Chinese hip-hop did not arrive on a straight line. In 1982, George Lam recorded a Cantonese freestyle on the television show \emph{Lam's Music Special}; in 1985, Cui Jian composed ``I Don't Understand,'' using rap-like delivery to handle city experience; in 1987, Harlem Yu's ``Report to Squad Leader'' gave hip-hop its first official Chinese name, ``xiha.''{\textsuperscript{[4]}} In 1992, Dai Bing and Tian Bao formed D.D. Rhythms, which later released \emph{Ku Re Jie Pai}, one of the earliest officially released mainland rap albums; in the 2000s, Longmenzhen, Yin Ts'ang, C-Block, and the Chengdu rap scene wrote dialect, region, and street experience into their songs.

In 2017, iQIYI's \emph{The Rap of China} pushed underground rap onto mainstream screens, and a generation of rappers reached a mass audience; GAI and Higher Brothers put Sichuan-Chongqing dialect, Chinese instruments, and trap rhythms into the same song, and Chinese hip-hop began to have its own accent.{\textsuperscript{[5]}} The route has not been smooth: commercialization, copyright, censorship, and community ethics keep pulling at it, but they also force a more concrete question---which parts of your beat could only come from your city?

\subsection{Sources for the History}
\label{sec:hiphop-sources}

\begin{longtable}{p{0.06\textwidth}p{0.38\textwidth}p{0.46\textwidth}}
\toprule
No. & Fact & Source \\
\midrule
\endfirsthead
\toprule
No. & Fact & Source \\
\midrule
\endhead
1 & August 11, 1973; 1520 Sedgwick Avenue; Kool Herc extends the break; Coke La Rock rhymes & GRAMMY, ``Founding Father DJ Kool Herc \& First Lady Cindy Campbell Celebrate Hip-Hop's 50th Anniversary,'' \url{https://www.grammy.com/news/dj-kool-herc-cindy-campbell-50th-anniversary-of-hip-hop-interview-jamaican-roots} \\
2 & 1990s golden age and representative works & Britannica, ``Hip-hop,'' \url{https://www.britannica.com/art/hip-hop} \\
3 & Trap's Atlanta route, T.I., and the producers & Rolling Stone, ``The 200 Greatest Rap Albums of All Time'' (Trap Muzik entry), \url{https://www.rollingstone.com/music/music-lists/best-hip-hop-albums-1323916/} \\
4 & Early attempts by George Lam, Cui Jian, and Harlem Yu; the formation of D.D. Rhythms & Apple Music, ``Hip-Hop Enters,'' \url{https://music.apple.com/cn/multi-room/6462871692}; QQ Music, ``Dai Bing artist profile,'' \url{https://y.qq.com/n/ryqq_v2/singer/000VwItL4LPm1k} \\
5 & 2017 \emph{The Rap of China} and mainstream exposure & The Paper, ``Developing `Zhou Yan,' the hidden line of `GAI','' \url{https://m.thepaper.cn/newsDetail_forward_31284981} \\
6 & Trap's 808, continuous hi-hats, and half-time feel & WNYC, ``Southern Rap and the Loss of an Icon,'' \url{https://www.wnycstudios.org/podcasts/bl/segments/southern-rap-and-loss-icon} \\
7 & \emph{Rapper's Delight}: production, Chic sample, mainstream breakthrough & GRAMMY, ``Essential Hip-Hop Releases From The 1970s,'' \url{https://www.grammy.com/news/hip-hop-albums-1970s-releases-videos-50th-anniversary}; Library of Congress, ``Rapper's Delight,'' \url{https://loc.gov/static/programs/national-recording-preservation-board/documents/RappersDelight.pdf} \\
8 & \emph{N.Y. State of Mind}: Joe Chambers's ``Mind Rain'' piano loop at the core & DJ Premier interview, ``DJ Premier Tells All: The Stories Behind His Classic Records,'' \url{https://djpremierblog.com/2011/02/19/dj-premier-tells-all-the-stories-behind-his-classic-records/} \\
9 & \emph{SICKO MODE}: three-part structure and two beat switches & Stereogum, ``The Number Ones: Travis Scott's `Sicko Mode','' \url{https://stereogum.com/2301905} \\
10 & GAI's \emph{Cang Hai Yi Sheng Xiao}: arrangement and production credits & Apple Music, ``Cang Hai Yi Sheng Xiao (Live),'' \url{https://music.apple.com/kr/song/1709991233} \\
11 & GAI's version reduces rap and adds melody and Chinese-style elements & The Paper, ``Developing `Zhou Yan,' the hidden line of `GAI','' \url{https://m.thepaper.cn/newsDetail_forward_31284981} \\
\bottomrule
\end{longtable}

\section{Arrangement: One Loop, Three Kinds of Space}
\label{sec:hiphop-craft}

\subsection{Write the Job Before the Instrument}
\label{sec:hiphop-instrumentation}

A hip-hop arrangement can be tiny: one voice, one drum machine, one sample. It can also be huge: orchestra, choir, 808s, and FX. Track count does not prove professionalism; every layer having a job does.

\begin{longtable}{p{0.16\textwidth}p{0.28\textwidth}p{0.26\textwidth}p{0.20\textwidth}}
\toprule
Layer & Common tracks/instruments & Main job & Common mistake \\
\midrule
\endfirsthead
\toprule
Layer & Common tracks/instruments & Main job & Common mistake \\
\midrule
\endhead
Voice & Lead, double, harmony, ad-lib & Tell the story, carry the tone, own the hook & Doubling every word, so it sounds like a copied person \\
Rhythm & Drum machine, 808, bass, percussion & Set tempo, set the pocket, give the body weight & Kick and 808 both full, so the low end has no hierarchy \\
Sample/harmony & Sample loop, keys, pad, guitar & Provide key, scene, and cycle & Stacking a pad on an already-full sample \\
Colour & Vinyl noise, ambience, strings, horns & Mark section borders and create hooks & Playing from start to finish, so the chorus has no arrival \\
\bottomrule
\end{longtable}

A workable start is six tracks: Vocal (audio), Double (audio), Drums (MIDI), 808 (MIDI), Sample (audio), and Master. Let those six stand up first; then decide whether to add a pad, strings, or a second sample. The more tracks you add, the more often you should solo-check: if this layer disappears, does the song get worse, or merely quieter? If you cannot answer, it can probably move to the background or go.

\subsection{Groove: Boom-Bap and Trap Are Two Pockets}
\label{sec:hiphop-groove}

\term{Boom-bap}{boom-bap} often puts the kick on 1 and 3 and the snare on 2 and 4, with eighth-note hi-hats and a little swing. Its tempo usually sits around 80--95 BPM and feels slow but biting. \term{Trap}{trap} often writes the tempo at 130--150 BPM but feels it in half time (roughly 65--75); long-decay 808s, continuous sixteenth or triplet hi-hats, sparse kicks, and sudden glides create a groove that rushes forward while sinking down.{\textsuperscript{[6]}}

The writing priorities differ:

\begin{itemize}
\item \textbf{Boom-bap}: keep the drums dry; do not crowd kick and snare onto the same beat; a sample loop can leave two beats of air for the voice.
\item \textbf{Trap}: hi-hat density is part of the rhythm language; do not write every bar with the same dense sixteenths; place 808 glides at phrase ends or section borders, not on every eighth.
\item \textbf{Both}: the kick is not four-on-the-floor. Four-on-the-floor is a dance/electronic habit; hip-hop needs the kick to be absent somewhere, and that absence is the vocal's pocket.
\end{itemize}

\term{Flow}{flow} is not speed; it is the relationship between syllables and drums. A syllable can land on the beat, anticipate it, lag behind it, compress, or stretch. When writing flow, count the bar first, then decide which words land on kick or snare; do not write every line with the same syllable count. The space at the end of a phrase is often more important than one more rhyme.

\subsection{Sampling: A Borrowed Loop Still Needs Your Gap}
\label{sec:hiphop-sample}

Sampling is not ``find a good song and drop it in.'' It has to pass three gates:

\begin{enumerate}
\item \textbf{Legal}: commercial release requires cleared samples or licensed material; even for practice, mark the source and do not treat someone else's full recording as your own work.
\item \textbf{Musical}: the loop needs a gap. If the original is full, EQ out low end, slice only a corner of the drums or voice, or leave half a beat open so new drums and 808 have somewhere to sit.
\item \textbf{Narrative}: why must this sample be here? Does it bring an era, a place, an accent, or only a ``vintage filter''? If it is only a filter, a clean self-recorded loop may serve better.
\end{enumerate}

\subsection{Harmony and Form: Simple Cycle, Clear Lift}
\label{sec:hiphop-harmony-form}

Hip-hop harmony can be as simple as i--iv--i, or just a bass line. The first job of harmony is not complexity; it is to give the loop direction and the voice a landing place. The 808 usually follows the root, but it need not glide on every eighth; too many glides turn the low end into a sentence without punctuation.

The most common form is Intro → Verse (16 bars) → Hook (8) → Verse (16) → Hook (8) → Bridge/Outro. Inside a verse, move in four-bar groups; do not change the sample every phrase. To make the song feel bigger, you do not necessarily add layers: you can remove one drum layer before the hook, or change the beat in the second verse so the listener sits up again.

\subsection{The Yeban Skeleton: Start With Six Tracks}
\label{sec:hiphop-yeban}

\begin{longtable}{p{0.12\textwidth}p{0.28\textwidth}p{0.48\textwidth}}
\toprule
Track & Type & What to write first \\
\midrule
\endfirsthead
\toprule
Track & Type & What to write first \\
\midrule
\endhead
Vocal & Audio & A sixteen-bar verse with half a beat of space at phrase ends; do not fix the tuning before the pocket works \\
Double & Audio & Double only hook keywords or phrase ends, never the whole verse \\
Drums & MIDI & Boom-bap: kick 1/3, snare 2/4, eighth hats; trap: sparse kick plus triplet hats \\
808 & MIDI & Follow the root and glide at phrase ends; solo it first and make sure it is not fighting the kick \\
Sample & Audio & A two- or four-bar loop; listen for gaps first; if there are none, shorten it or cut low end \\
Master & Bus & Check level last; do not use the bus to hide a low-end conflict during writing \\
\bottomrule
\end{longtable}

Order: place the loop, write the drums, write the 808, then write the voice. Solo after every layer. If removing a track only makes the song quieter, it was probably a habit, not an arrangement decision.

\section{Four Songs, Four Arrangement Lessons}
\label{sec:hiphop-cases}

\subsection{\emph{Rapper's Delight}: Putting a Party on a Record}
\label{sec:hiphop-case-rappers-delight}

In 1979, the Sugarhill Gang's \emph{Rapper's Delight}, produced by Sylvia Robinson, used Chic's ``Good Times'' as a dance-floor loop while three MCs traded lines, pressing the heat of a live party into the studio.{\textsuperscript{[7]}} It was not the first recorded rap song, but it was the first single that let a mass audience meet rap in a record shop.

The arrangement's intelligence is restraint: drums and bass keep the disco cycle, while voices trade long lines, short lines, and call-and-response; sections change through MC rotation, not through piling on instruments. It sounds like a party that refuses to end because the loop never tries to prove how complex it is.

\textbf{One thing to steal}: if you have only one loop, write who enters when. Give the three voices different phrase lengths and entrance points; do not give all three the same flow.

\subsection{\emph{N.Y. State of Mind}: The Dark Space Inside One Loop}
\label{sec:hiphop-case-ny-state}

Nas's 1994 \emph{N.Y. State of Mind} was produced by DJ Premier. The beat centers on a piano loop from Joe Chambers's ``Mind Rain,'' with dry drums and scratches.{\textsuperscript{[8]}} It does not explain New York with a large ensemble; it lets one minor-key piano loop repeat and leaves the city's pressure in the gaps.

The drums are dry, the piano loop does not fill every beat, and Nas's voice sometimes anticipates and sometimes lags, refusing to rhyme on every phrase end. What the listener hears is not ``many layers'' but the tension between one loop and one voice.

\textbf{One thing to steal}: subtract first. Remove the pad, the second drum layer, and the extra sample from your demo one by one, until only loop, drums, and voice remain. If that still stands, then decide what to add.

\subsection{\emph{SICKO MODE}: Reallocating Attention With Beat Switches}
\label{sec:hiphop-case-sicko}

Travis Scott's 2018 \emph{SICKO MODE} is not an even three minutes. It has at least three clear sections, with the beat switching twice: a slow Drake entrance, a harder trap drum section, then a darker, lower beat.{\textsuperscript{[9]}} Samples, Auto-Tune, ad-libs, and ambience are not decoration; they do the work of breathing and changing scene.

Its arrangement logic is not ``make the chorus longer'' but ``split one song into three short songs,'' each solving one emotional task. The beat switch forcibly resets attention, so each section can use a different density, low end, and vocal treatment.

\textbf{One thing to steal}: a beat switch is not ``change to another loop for no reason.'' The before and after must differ in at least two things---sense of speed, low-end weight, vocal treatment---or the listener only hears a stutter.

\subsection{GAI's \emph{Cang Hai Yi Sheng Xiao}: Jianghu Inside Trap}
\label{sec:hiphop-case-gai}

After \emph{The Rap of China} in 2017, GAI turned \emph{Cang Hai Yi Sheng Xiao} into a version built for a mainstream stage. The original comes from the 1990 film \emph{The Swordsman} and is a pentatonic melody written by James Wong; GAI's version keeps that skeleton while adding Sichuan-Chongqing dialect rap, pipa/percussion, and a modern drum kit. The arrangement is by Song Yang, with Liu Zhou as producer.{\textsuperscript{[10]}} The Paper notes that this version deliberately reduced the more aggressive rap section and added melody and Chinese-style elements to fit a mainstream stage.{\textsuperscript{[11]}}

What works is not a ``ancient-style filter'' but two times coexisting in one song: the pentatonic melody is old, the 808/trap low end is new, and the dialect pronunciation turns ``jianghu'' from an abstract word into a specific local accent.

\textbf{One thing to steal}: when fusing, do not borrow two complete languages at once. Borrow one skeleton (say, a pentatonic melody) and use a groove you know for everything else; otherwise it becomes a pile of elements that resembles everything and owns nothing.

\section{Three Exercises at Yeban}
\label{sec:hiphop-practice}

\begin{exercise}
\textbf{Eight bars, three starting points.} Make three versions of the same eight bars: A starts from a licensable drum/ambience sample; B writes drums and 808 in MIDI first; C hums a flow and reverse-engineers the drums. Compare recognition, speed, editability, and copyright risk. A is fast but has clearance cost; B is slow but every beat can change; C sounds most like you but easily becomes a loop with no space.
\end{exercise}

\begin{exercise}
\textbf{One verse, three flows.} Write eight lines and perform them three ways: A mostly on the beat with space at phrase ends; B heavy anticipation/lag to create a rush; C mostly triplets with compressed phrases. Record all three, have someone listen blind, and ask which sounds most like ``talking'' and which most like ``singing.'' There is no right answer, only trade-offs.
\end{exercise}

\begin{exercise}
\textbf{Two beat switches, no lost listener.} Build twenty-four bars: 0--8 boom-bap, 9--16 trap, 17--24 back to boom-bap but with a different sample. Each switch must change at least two of low end, drum density, and vocal treatment. Listen back; if the second switch feels like ``returning to the start,'' change one more variable.
\end{exercise}

\section{Further Reading}
\label{sec:hiphop-reading}

\begin{sloppypar}
\begin{itemize}
\item \textbf{Rolling Stone, ``The 200 Greatest Rap Albums of All Time''}: \url{https://www.rollingstone.com/music/music-lists/best-hip-hop-albums-1323916/}. Read the entries by decade to fill in works, scenes, and producers.
\item \textbf{GRAMMY, ``Founding Father DJ Kool Herc \& First Lady Cindy Campbell Celebrate Hip-Hop's 50th Anniversary''}: \url{https://www.grammy.com/news/dj-kool-herc-cindy-campbell-50th-anniversary-of-hip-hop-interview-jamaican-roots}. Check the 1973 party, breakbeat, and early MC details against the participants' own interview.
\item \textbf{Music Theory Online, ``We Gon' Be Alright?''}: \url{https://mtosmt.org/issues/mto.19.25.1/mto.19.25.1.manabe.html}. Read an academic analysis that puts hook, flow, harmony, and reception history together.
\item \textbf{Library of Congress, ``Rapper's Delight''}: \url{https://loc.gov/static/programs/national-recording-preservation-board/documents/RappersDelight.pdf}. Use the archive essay to check the single's recording history and mainstream breakthrough.
\end{itemize}
\end{sloppypar}

These links were checked against their official pages on November 27, 2026. If a regional restriction blocks access, search the item title in the platform's official account.

\section{Chapter Takeaways}
\label{sec:hiphop-takeaways}

\begin{itemize}
\item Hip-hop is not ``fast talking''; it is a culture grown by DJs, MCs, dancers, and writers, and rap is only its most recordable line.
\item A loop's professionalism comes from its gaps: the drums must be absent somewhere, the sample must leave air, and the voice must have a pocket.
\item Boom-bap and trap are two pockets, not two preset packs; decide the groove before the instruments.
\item Sampling passes through legal, musical, and narrative gates; commercial release needs clearance, and practice still needs source credit.
\item Form need not be complex, but its lift must be clear; a beat switch changes at least two things, or it is only a stutter.
\item When fusing, borrow one skeleton, not two complete languages.
\end{itemize}

\begin{helpmapping}
\begin{longtable}{p{0.45\textwidth}p{0.45\textwidth}}
\toprule
Section & Help anchor \\
\midrule
What hip-hop is & \texttt{arrangement-genre-hiphop-what} \\
The route & \texttt{arrangement-genre-hiphop-history} \\
Arrangement & \texttt{arrangement-genre-hiphop-craft} \\
Four songs & \texttt{arrangement-genre-hiphop-cases} \\
Yeban skeleton & \texttt{arrangement-genre-hiphop-yeban} \\
Exercises & \texttt{arrangement-genre-hiphop-practice} \\
Further reading & \texttt{arrangement-genre-hiphop-reading} \\
\bottomrule
\end{longtable}
\end{helpmapping}

\begin{figure}[htbp]\centering
\begin{tikzpicture}[font=\scriptsize,x=1cm,y=1cm,>=Stealth]
\yebangenreddraw{Hip-hop: let the low end speak first}{20--80 Hz 808 sub}{200--2 kHz voice and samples}{2--10 kHz hats and detail}
\end{tikzpicture}
\caption{Hip-hop frequency split, drum grid and section-energy curve}
\end{figure}
